%% kapa: exam=2 round=35 year=2024 period=3 subject=compensation_law subject_ko=감정평가 및 보상법규 seq=4 number=4 points=10
%% source: Q-Net 감정평가사 기출문제 (cst00309) artlSeq=5213532 file=raw/1008/5213532/2024년 제35회 감정평가사 2차시험 문제지.pdf
%% conversion: kordoc 4.24.51 -> tools/gen2.py (auto); review=proofread-sonnet; flags=none
\begin{longquestion}{【문제 4】}
「감정평가 및 감정평가사에 관한 법률」 제28조 제1항에 따른 손해배상책임을 보장하기 위하여 감정평가법인등이 하여야 하는 ‘필요한 조치’의 내용과 ‘필요한 조치’를 하지 아니한 경우 「감정평가 및 감정평가사에 관한 법률」에 따른 행정상 제재를 설명하시오. (10점) \par
\medskip
\begin{framedtext}
\textbf{참조 조문} \par
\textbf{<감정평가 및 감정평가사에 관한 법률>} \par
\textbf{제28조(손해배상책임)} ① 감정평가법인등이 감정평가를 하면서 고의 또는 과실로 감정평가 당시의 적정가격과 현저한 차이가 있게 감정평가를 하거나 감정평가 서류에 거짓을 기록함으로써 감정평가 의뢰인이나 선의의 제3자에게 손해를 발생하게 하였을 때에는 감정평가법인등은 그 손해를 배상할 책임이 있다. \par
\end{framedtext}
\end{longquestion}
